% `tex.web` §863 and §873: which passes a paragraph announces.
%
%     if pretolerance>=0 then
%       begin if tracing_paragraphs>0 then
%         begin begin_diagnostic; print_nl("@firstpass"); end;
%       threshold:=pretolerance; second_pass:=false; final_pass:=false;
%       end
%     else begin threshold:=tolerance; second_pass:=true;
%       final_pass:=(emergency_stretch<=0);
%       if tracing_paragraphs>0 then begin_diagnostic;
%       end;
%
% Only the first pass is announced there, and only if there is one: with
% `\pretolerance` negative the narration starts straight into the breaks with
% no header at all. `@secondpass` is §873's and is printed when a first pass
% has *failed*, so a paragraph that never had one cannot produce it — we
% printed it anyway.
%
% §873 also decides whether there is a third pass:
%
%     final_pass:=(emergency_stretch<=0);
%     … else begin if tracing_paragraphs>0 then print_nl("@emergencypass");
%       background[2]:=background[2]+emergency_stretch; final_pass:=true; end;
%
% With `\emergencystretch` at zero — INITEX's value, and trip.tex's — the
% second pass *is* the final one and `@emergencypass` never appears. We ran a
% third pass unconditionally and labelled it, so every paragraph that needed
% hyphenation claimed an emergency.
%
% What that third pass *does* is unchanged and still not TeX's: §851's
% artificial demerits make a forced break free inside the final pass, and we
% score it. That shows as `b=*`/`d=*` in Knuth's transcript against real
% numbers here, and sometimes as a different break. This case stays clear of
% that historical final-pass defect: both paragraphs below are broken by a
% pass that succeeds.
\catcode`\{=1 \catcode`\}=2
\tracingonline=1
\tracingparagraphs=1
\hsize=100pt \parindent=0pt
\font\rip=cmr10 \rip

% A first pass that succeeds: announced, and nothing else.
\pretolerance=10000
AAA AAA\par

% No first pass at all: no header.
\pretolerance=-1 \tolerance=10000
AAA AAA\par
\end
